\scanpage{321}
\mspage{\deleted{\unsure{270}}271}

$i_1,i_2$ sont des \uline{imm. ouvertes} (de Dwyer)\inserted{\deleted{\unsure{asphériques|symétriques}}} et où le carré est cocartésien (ce qui implique qu'il est $W$-cocartésien). Faisant alors le \unsure{cochgt.} de base par $Z_0\xrightarrow{p_0}Z'_0=e$, on trouve un \deleted{carré} hom de carrés cocartésiens

\diagram{\(\begin{tikzcd}[ampersand replacement=\&,column sep=small,row sep=small]
{}\&Z_1\arrow[dr]\arrow[ddd,dashed,"{p_1}" near start]\&{}\\
Z_0\arrow[ur]\arrow[dr]\arrow[ddd,"{\deleted{\unsure{p}}p_0}"']\&{}\&Z_3\arrow[ddd,"{p_3}"]\\
{}\&Z_2\arrow[ur]\arrow[ddd,dash,"{p_2}" near end]\&{}\\
{}\&Z'_1\arrow[dr,dashed]\&{}\\
e=Z'_0\arrow[ur]\arrow[dr]\&{}\&Z'_3\\
{}\&Z'_2\arrow[ur]\&{}
\end{tikzcd}\)}

où toutes les faces du cube sont $W$-cocartésiennes, et comme $p_0\in W\unsure{,}$ ($Z_0$ \unsure{étant} $W$-asph) les autres flèches verticales $p_1,p_2,p_3$ sont $\in W$, donc le carré $Z'_{\unsure{*|\bullet}}$ satisfait aussi $Z'_0$ et $Z'_3$ asph (et de plus $Z'_0=e$), et il faut prouver $Z'_1,Z'_2$ $W$-asph (d'où idem pour $Z_1,Z_2$). Donc, on est ramené au cas particulier où $Z_0=e$.

Mais dans ce cas il existe une rétraction \inserted{(triviale !)} de $Z_1$ sur $Z_0$, donc aussi de $Z_3$ sur $Z_2$. Ainsi, $Z_2$ est un facteur
% archive sheet number: 321
